% Generated by render.py; edit preamble.tex, citations.json and routes.tikz
\documentclass[border=12pt]{standalone}
\def\pgfsysdriver{pgfsys-dvisvgm.def}
\usepackage[dvisvgm]{xcolor}
\usepackage{xeCJK}
\setCJKmainfont{HaranoAjiGothic-Regular.otf}[BoldFont=HaranoAjiGothic-Bold.otf]
\setCJKsansfont{HaranoAjiGothic-Regular.otf}[BoldFont=HaranoAjiGothic-Bold.otf]
\usepackage{amsmath,amssymb,tikz,lmodern}
\usetikzlibrary{arrows.meta,positioning,calc,fit,backgrounds}
\renewcommand{\familydefault}{\sfdefault}
\definecolor{ink}{HTML}{25322F}
\definecolor{green}{HTML}{246657}
\definecolor{tint}{HTML}{EEF4F0}
\definecolor{soft}{HTML}{C6D2CC}
\definecolor{muted}{HTML}{52615B}
\newif\ifJapanese
\newcommand{\Tx}[2]{\ifJapanese #1\else #2\fi}
\newcommand{\Eff}{\overline{\operatorname{Eff}}}
\newcommand{\heading}[3]{{\color{green}\bfseries #1\quad\Tx{#2}{#3}}\\[5pt]}
% A white rectangle makes the complete citation label clickable, including spaces.
\newcommand{\refurl}[2]{\special{dvisvgm:raw <a xlink:href="#1" target="_blank" rel="noopener noreferrer">}{\setlength{\fboxsep}{0pt}\colorbox{white}{\color{green}\underline{\strut #2}}}\special{dvisvgm:raw </a>}}
\tikzset{
 card/.style={draw=soft,fill=white,rounded corners=2pt,align=left,text=ink,font=\fontsize{10.5}{14}\selectfont,inner xsep=12pt,inner ysep=10pt,text width=14.4cm},
 half/.style={card,text width=6.6cm,minimum height=2.05cm},
 result/.style={card,draw=green,fill=tint,line width=.9pt},
 reason/.style={align=left,text=muted,font=\fontsize{9}{12}\selectfont,inner sep=0pt,text width=12.6cm},
 branchreason/.style={reason,text width=6.8cm},
 edge/.style={-{Stealth[length=2.4mm]},draw=green,line width=1pt},
 wire/.style={draw=green,line width=1pt},
}
% Keep the arrow in a dedicated left gutter; explanations never cover it.
\def\flow#1#2#3;{
 \coordinate (#2-from) at ($(#1.south)+(-6.9,0)$);
 \coordinate (#2-to) at ($(#2.north)+(-6.9,0)$);
 \draw[edge] (#2-from)--(#2-to);
 \node[reason,anchor=west] (#2-reason) at ($(#2-from)!.5!(#2-to)+(.55,0)$) {#3};
}

% Style adapted from the repository's Schnell diagrams; MIT attribution retained.

% Generated from citations.json; do not edit.
\newcommand{\Cite}[1]{\csname cite-#1\endcsname}
\expandafter\def\csname cite-three-main\endcsname{\shortstack[l]{\refurl{https://github.com/openai/math/blob/1581d4cbdd69a0faf34b2089614fe2a814316a6a/preprints/The-Kakeya-maximal-conjecture-in-three-dimensions-September-23-2026/paper.pdf}{[OpenAI26a, Theorem 1.1, p. 2]}}}
\expandafter\def\csname cite-three-plank\endcsname{\shortstack[l]{\refurl{https://github.com/openai/math/blob/1581d4cbdd69a0faf34b2089614fe2a814316a6a/preprints/The-Kakeya-maximal-conjecture-in-three-dimensions-September-23-2026/paper.pdf}{[OpenAI26a, Lemma 3.3, pp. 20–21]}}}
\expandafter\def\csname cite-three-use\endcsname{\shortstack[l]{\refurl{https://github.com/openai/math/blob/1581d4cbdd69a0faf34b2089614fe2a814316a6a/preprints/The-Kakeya-maximal-conjecture-in-three-dimensions-September-23-2026/paper.pdf}{[OpenAI26a, Lemmas 4.7, 4.9, pp. 32–33, 35–38]}}}
\expandafter\def\csname cite-three-end\endcsname{\shortstack[l]{\refurl{https://github.com/openai/math/blob/1581d4cbdd69a0faf34b2089614fe2a814316a6a/preprints/The-Kakeya-maximal-conjecture-in-three-dimensions-September-23-2026/paper.pdf}{[OpenAI26a, Section 10, pp. 86–94]}}}
\expandafter\def\csname cite-four-main\endcsname{\shortstack[l]{\refurl{https://github.com/openai/math/blob/1581d4cbdd69a0faf34b2089614fe2a814316a6a/preprints/Every-four-dimensional-Kakeya-set-has-full-Hausdorff-dimension-September-24-2026/paper.pdf}{[OpenAI26b, Theorem 1.1, p. 3]}}}
\expandafter\def\csname cite-four-plank\endcsname{\shortstack[l]{\refurl{https://github.com/openai/math/blob/1581d4cbdd69a0faf34b2089614fe2a814316a6a/preprints/Every-four-dimensional-Kakeya-set-has-full-Hausdorff-dimension-September-24-2026/paper.pdf}{[OpenAI26b, Lemma 2.15, pp. 17–18]}}}
\expandafter\def\csname cite-four-reduction\endcsname{\shortstack[l]{\refurl{https://github.com/openai/math/blob/1581d4cbdd69a0faf34b2089614fe2a814316a6a/preprints/Every-four-dimensional-Kakeya-set-has-full-Hausdorff-dimension-September-24-2026/paper.pdf}{[OpenAI26b, Theorem 2.17, pp. 20–25]}}}
\expandafter\def\csname cite-four-end\endcsname{\shortstack[l]{\refurl{https://github.com/openai/math/blob/1581d4cbdd69a0faf34b2089614fe2a814316a6a/preprints/Every-four-dimensional-Kakeya-set-has-full-Hausdorff-dimension-September-24-2026/paper.pdf}{[OpenAI26b, Section 10.8, p. 168]}}}
\expandafter\def\csname cite-catalogue\endcsname{\shortstack[l]{\refurl{https://github.com/openai/math/blob/1581d4cbdd69a0faf34b2089614fe2a814316a6a/overview.pdf}{OpenAI Research Catalog · 074 · p. 9}}}

\Japanesetrue
\begin{document}
\begin{tikzpicture}
\node[result] (input) at (0,0) {
\heading{}{全時間plankへの集中を制御する}{Control concentration in full-time planks}
$M_i(t)=b_i+tu_i,\quad \sum_{M_i\subset P}\omega_i\le Aab,\quad \#J_i\ge D^{1-\eta}$\\[4pt]
$\displaystyle\sum_i\omega_i\#J_i\le C_\varepsilon A D^\varepsilon\#E$
\quad \Cite{three-plank}
};
\node[half] (three) at (-4.15,-5.4) {
\heading{}{三次元の配置へ使う}{Three-dimensional input}
\Tx{flat scaleのplate捕捉と\newline 正のdeficitの配置}{Plate capture at a flat scale;\newline a positive-deficit configuration}\\[5pt]
\Tx{別の幾何入力と定常配置の排除を加える}{Additional geometric inputs and exclusion of stationary configurations}
};
\node[half] (four) at (4.15,-5.4) {
\heading{}{四次元の反例を移す}{Four-dimensional reduction}
$\dim_H K<4\ \Longrightarrow\ H(\mathcal E)>0$\\[5pt]
\Tx{可変尺度被覆・有限選択・重みを保持}{Retain variable-scale covers, finite selections, and weights}
};
\draw[edge] ($(input.south)+(-7.5,0)$)--($(three.north)+(-3.35,0)$);
\draw[edge] ($(input.south)+(.8,0)$)--($(four.north)+(-3.35,0)$);
\node[branchreason,anchor=west] at (-7.05,-2.7) {\Tx{全時間の集中評価を使う}{Use the full-time concentration estimate}\\[4pt]\Cite{three-use}};
\node[branchreason,anchor=west] at (1.25,-2.7) {\Tx{入力の内容対応を照合}{Matched input statement}\\[4pt]\Cite{four-plank}\\[4pt]\Tx{利用先はLemma 2.3と引用：番号不整合}{Receiving citation says Lemma 2.3: number mismatch}};
\node[half,draw=green,fill=tint] (r3) at (-4.15,-11.5) {
\heading{}{三次元：最大関数評価}{3D: maximal estimate}
$\|K_\delta f\|_{L^3(S^2)}\le C_\varepsilon\delta^{-\varepsilon}\|f\|_{L^3(\mathbb R^3)}$\\[5pt]
\Tx{すべての $\varepsilon>0$、$0<\delta<1$}{Every $\varepsilon>0$ and $0<\delta<1$}\\[4pt]\Cite{three-main}
};
\node[half,draw=green,fill=tint] (r4) at (4.15,-11.5) {
\heading{}{四次元：Hausdorff次元}{4D: Hausdorff dimension}
$\dim_H K=4$\\[5pt]
\Tx{全方向に単位線分を含む任意集合}{An arbitrary set containing a unit segment in every direction}\\[4pt]\Cite{four-main}
};
\draw[edge] ($(three.south)+(-3.35,0)$)--($(r3.north)+(-3.35,0)$);
\draw[edge] ($(four.south)+(-3.35,0)$)--($(r4.north)+(-3.35,0)$);
\node[branchreason,anchor=west] at (-7.05,-8.3) {\Tx{臨界指数を消し、一般の $L^3$ 関数へ}{Remove the critical loss; pass to general $L^3$ functions}\\[4pt]\Cite{three-end}};
\node[branchreason,anchor=west] at (1.25,-8.3) {\Tx{正のhorizonを三つの配置の改善で排除}{Exclude positive horizon through three improvement regimes}\\[4pt]\Cite{four-reduction}\\[4pt]\Cite{four-end}};

\end{tikzpicture}
\end{document}
